\section*{Capítulo \ref{ch:b2:thermodynamic-potentials} --- Potenciais termodinâmicos}

\begin{solution}{exo:b2:thermodynamic-potentials:1}
$F = -nRT\ln V + f(T)$, $G = nRT\ln P + g(T)$; $P = -\partial_VF = nRT/V$,
$V = \partial_PG = nRT/P$; $-\partial_TF = nR\ln V - f'(T)$, a entropia
$nC_V\ln T + nR\ln V + \text{const.}$ se $f' = -nC_V\ln T + \dots$
\end{solution}

\begin{solution}{exo:b2:thermodynamic-potentials:2}
$\Delta F = -RT\ln 10 = \qty{-5.7}{kJ}$, $\Delta U = 0$, $\Delta S = \qty{19.1}{J/K}$.
Reversível: \qty{5.7}{kJ}, $S_{\text{c}} = 0$; contra \qty{1}{bar}: \qty{0.9}{kJ},
$S_{\text{c}} = 19.1 - 900/300 = \qty{16}{J/K}$; no vácuo: nada, $S_{\text{c}} =
\qty{19.1}{J/K}$.
\end{solution}

\begin{solution}{exo:b2:thermodynamic-potentials:3}
$\Delta g = \Delta h - T\Delta s$ com $\Delta h \approx L$, $\Delta s \approx L/T_{\text{fus}}$.
A \qty{-5}{\celsius}: $g_{\text{liq}} - g_{\text{gelo}} = +\qty{6.1}{kJ/kg}$, o gelo é
estável e o líquido metastável (o super-resfriamento existe); a $+\qty{5}{\celsius}$:
\qty{-6.1}{kJ/kg}, o líquido é estável --- gelo superaquecido não se observa, pois sua
superfície funde primeiro.
\end{solution}

\begin{solution}{exo:b2:thermodynamic-potentials:4}
Fusão $\qty{-1.3e7}{Pa/K}$; vaporização \qty{3.6}{kPa/K}. \qty{70}{bar}:
\qty{-0.5}{K}; \qty{0.7}{bar}: \qty{92}{\celsius}; \qty{2}{bar}: $\approx
\qty{127}{\celsius}$ pela inclinação local (\qty{120}{\celsius} de fato: a
inclinação cresce).
\end{solution}

\begin{solution}{exo:b2:thermodynamic-potentials:5}
(a) Também $(\partial T/\partial V)_S = -(\partial P/\partial S)_V$ e
$(\partial T/\partial P)_S = (\partial V/\partial S)_P$. (b) $0$; $a/v^2$. (c)
$\alpha/\kappa_T = \qty{4.6}{bar/K}$; $T\alpha/\kappa_T - P \approx \qty{1400}{bar}$.
(d) \qty{46}{bar}: a garrafa estoura.
\end{solution}

\begin{solution}{exo:b2:thermodynamic-potentials:6}
(a) \cref{thm:b2:thermodynamic-potentials:maxwell}. (b) $nR$. (c) Água
\qty{29}{J/kg/K} ($0.7\%$); cobre \qty{12}{J/kg/K} ($3\%$). (d) A
diferença é desprezível para a matéria condensada.
\end{solution}

\begin{solution}{exo:b2:thermodynamic-potentials:7}
(a) $-C(L - L_0)$; $0$. (b) $\Delta S = -C(L - L_0)^2/2 = \qty{-5e-3}{J/K}$; $Q =
\qty{-1.5}{J}$; $W = +\qty{1.5}{J}$. (c) $+\qty{0.4}{K}$. (d) $(L - L_0) \propto
1/T$: $-9\%$.
\end{solution}

\begin{solution}{exo:b2:thermodynamic-potentials:8}
(a) A matéria vai para o $\mu$ menor: $\mu_{\text{v}} < \mu_{\text{liq}}$ quando $p <
p_{\text{sat}}$; umidade $= p/p_{\text{sat}}$. (b) $-RT\ln 0.5 = \qty{1.7}{kJ/mol}$
em favor do vapor. (c) A evaporação só exige $p < p_{\text{sat}}$,
e não a ebulição; a $100\%$ os potenciais são iguais. (d) Nas superfícies
mais frias, cuja $p_{\text{sat}}(T_{\text{s}})$ caiu abaixo do $p$ do ar
--- aquelas que irradiam para o céu noturno.
\end{solution}

\begin{solution}{exo:b2:thermodynamic-potentials:9}
(a) Padrão. (b) $b = RT_{\text{c}}/8P_{\text{c}} = \qty{4.3e-5}{m^3/mol}$, $a =
27R^2T_{\text{c}}^2/64P_{\text{c}} = \qty{0.36}{Pa.m^6/mol^2}$. (c) $\dd g = v\,\dd P$
a $T$ fixo e $g$ igual nas duas extremidades: $\int v\,\dd P = 0$ ao longo do
laço. (d) Instável onde $\partial_vP > 0$; metastável entre os dois: o
líquido superaquecido (que explode) e o vapor supersaturado (a névoa que
espera por núcleos, a câmara de nuvens).
\end{solution}

\begin{solution}{exo:b2:thermodynamic-potentials:10}
(a) $\dd P/\dd T = L_{\text{m}}P/RT^2$: $\ln(P/P_0) = -(L_{\text{m}}/R)(1/T - 1/T_0)$.
(b) \qty{28}{mbar} contra \qty{23}{mbar} --- $L$ é maior a \qty{20}{\celsius}.
(c) \qty{820}{Pa} contra \qty{611}{Pa}. (d) $\partial(\Delta g/T)/\partial T = -L/T^2$
e $\Delta g = 0$ ao longo da curva dão a mesma inclinação.
\end{solution}

\begin{solution}{exo:b2:thermodynamic-potentials:11}
(a) $\mu$ igual dos dois lados da membrana: $\mu^\circ(P + \Pi) - RTx_{\text{s}} =
\mu^\circ(P)$, ou seja, $v_{\text{m}}\Pi = RTx_{\text{s}}$, $\Pi = cRT$. (b) \qty{27}{bar}.
(c) $\Pi V = \qty{2.7}{MJ} = \qty{0.75}{kWh}$ por metro cúbico; as instalações reais
gastam quatro vezes mais. (d) A água escoa para o potencial menor: para dentro da
célula rica em soluto, e para fora dela na salmoura.
\end{solution}

\begin{solution}{exo:b2:thermodynamic-potentials:12}
(a) $S = S_0 - k_BL^2/2Na^2$; $f = k_BTL/Na^2$. (b) $k_BT/Na^2 \propto T$. (c)
\qty{1.7e-13}{N}; $1.8 \times 10^{14}$. (d) $E \sim nk_BT \approx \qty{0.4}{MPa}$.
\end{solution}

\begin{solution}{pb:b2:thermodynamic-potentials:1}
\textbf{1.} $\dd U = T\dd S + f\dd L$; $\dd F = -S\dd T + f\dd L$.

\textbf{2.} Schwarz aplicado a $\dd F$.

\textbf{3.} $f - T\partial_Tf = CT(L - L_0) - CT(L - L_0) = 0$: a energia
não muda; a força de retorno é entrópica.

\textbf{4.} $\Delta S = \qty{-5e-3}{J/K}$; $Q = \qty{-1.5}{J}$; $W = +\qty{1.5}{J}$; soma
nula.

\textbf{5.} $+\qty{0.4}{K}$.

\textbf{6.} $L - L_0 \propto 1/T$: \qty{9.1}{cm}, uma contração de \qty{9}{mm}; os
raios aquecidos encurtam, o centro de massa do aro se desloca e a roda
gira.

\textbf{7.} $S = S_0 - k_BL^2/2Na^2$; $f = k_BTL/Na^2$; a penalidade de entropia é
paga em $k_BT$.

\textbf{8.} \qty{1.7e-13}{N}; $1.8 \times 10^{14}$ cadeias.

\textbf{9.} Sua força é energética ($\partial_LU \ne 0$); aquecer amolece
as ligações, de modo que $f$ cai com $T$.

\textbf{10.} Inclinações $-s$, com $s_{\text{gas}} > s_{\text{liq}} > s_{\text{gelo}}$; a
linha mais baixa é a fase estável; cruzamentos em $T_{\text{fus}}$ e $T_{\text{vap}}$.

\textbf{11.} \cref{thm:b2:thermodynamic-potentials:coexistence}.

\textbf{12.} $\qty{-1.3e7}{Pa/K}$; negativa porque $v_{\text{liq}} < v_{\text{gelo}}$.

\textbf{13.} \qty{70}{bar}: \qty{-0.5}{K} --- não é assim que se patina a
\qty{-10}{\celsius}; o atrito e a camada superficial quase líquida é que fazem isso.

\textbf{14.} \qty{270}{bar}, \qty{-2}{K}: com o fluxo geotérmico, a
base funde e o manto desliza.

\textbf{15.} \qty{3.6}{kPa/K}; \qty{92}{\celsius}; cerca de \qty{127}{\celsius}
(\qty{120}{} de fato).

\textbf{16.} \qty{28}{mbar} (\qty{23}{}), \qty{820}{Pa} (\qty{611}{}): $L$ cresce quando $T$
cai; erros de $20$--$30\%$.

\textbf{17.} O mesmo $\Delta v$, de modo que a razão das inclinações vale $L_{\text{sub}}/L_{\text{vap}}
= 1.15$: a curva de sublimação é mais íngreme --- um bico no ponto triplo.

\textbf{18.} \qty{12}{kJ/kg}; o gelo precisa nuclear, e um núcleo pequeno custa
mais energia de superfície do que ganha: as gotículas puras esperam, até
\qty{-40}{\celsius}.

\textbf{19.} $\Delta v \to 0$, $L \to 0$, com razão finita: a curva termina;
além dele, um único fluido que passa continuamente de um estado tipo líquido a um tipo gás.

\textbf{20.} $\dd g = v\,\dd P$ a $T$ fixo; $g$ igual nas duas extremidades da
horizontal: $\int v\,\dd P = 0$, áreas iguais.

\textbf{21.} O ramo do líquido superaquecido: nenhum núcleo de bolha numa
xícara lisa; uma perturbação o faz ferver de uma vez.

\textbf{22.} Congelar, bombear abaixo de \qty{611}{Pa}, aquecer suavemente: o caminho
cruza a curva de sublimação, nunca a de fusão; a estrutura
é preservada.

\textbf{23.} $\mu_{\text{v}}(p) < \mu_{\text{liq}}$ quando $p < p_{\text{sat}}$; iguais na
saturação.

\textbf{24.} \qty{120}{\celsius} em vez de \qty{100}{}: quatro vezes mais depressa.

\textbf{25.} $S$ para o isolado, $F$ para $T,V$ fixos, $G$ para $T,P$ fixos;
$-\Delta F$, $-\Delta G$ são o trabalho recuperável; as relações de Maxwell transformam
derivadas de $S$ na equação de estado.
\end{solution}
